% =====================================================================
%  GMN2 §12 范畴的若干问题 + §13 $R\to\infty$ 极限 + §14 与 [1] 的比较
%  原文：wkb-paper.tex 8421–9207 行；PDF（印刷）p.152–163；图 84
%  编号公式：(12.1) 1 条 + (13.1)–(13.46) 46 条 + (14.1)–(14.8) 8 条，共 55 条
%  原文宏 \nsing（= l）preamble 未定义，此处本地补齐
% =====================================================================

\providecommand{\nsing}{l}

\section{范畴的若干问题} \label{sec:category}

偏数学的读者也许会问：本文的构造与 \cite{ks1} 中处理 Donaldson–Thomas 不变量的进路究竟能接近到什么程度？特别是，那条进路以某个适当的范畴及其上的一族稳定性条件为出发点。这些数据在我们的例子里如何实现？本节提出若干观察，或可有助于指明方向。我们的讨论非常粗略，并回避了几个重要问题，尤其是 $\IZ$-分级在这套故事中的作用。

先从 $K=2$ 的情形说起（最近几节我们一直集中在这个情形）。类比于 Calabi–Yau 三维流形的 Fukaya 范畴，我们可以把基本故事搭成如下样子。我们有一个黎曼面 $C$，以及 $C$ 上亚纯二次微分构成的空间 $\CB$。令 $\CB'\subset\CB$ 为正则区域，由只含\emph{单}零点的二次微分组成。对任意 $u\in\CB'$，有相应的双覆盖 $\Sigma_u\to C$，并带一个 1-形式 $\lambda$。在通常关于 Fukaya 范畴上稳定性条件的讨论中，人们考虑一族在辛意义下同构的 Calabi–Yau 流形，因而局部可以把它们与一个固定的辛流形等同起来。这里类似地，各个覆盖 $\Sigma_u$ 若只当作 $C$ 的拓扑分歧覆盖来看，也可以与一个固定的 $\Sigma$ 等同。于是：

\begin{itemize}

\item 我们范畴的对象，应当是 $\Sigma$ 上一族互不相交的有向闭路 $\gamma$ 的汇总（或者也许是这类路汇的复形），它们在交换两叶之下反变。

\item 两条路 $\gamma_1$、$\gamma_2$ 之间的态射空间，应当对 $\gamma_1\cap\gamma_2$ 的每个点各有一个基元素。

\item 一条路 $\gamma$ 的 K 理论类应当是它的同调类。

\item 稳定性条件的空间应当是 $\CB'$ 的万有覆叠 $\tilde{\CB}$。

\end{itemize}

给定 $\tilde{\CB}$ 的一点，我们可以把路 $\gamma$ 上的\emph{相位函数}定义为 $\lambda\cdot\partial_t$ 的相位，其中 $\partial_t$ 是正向切向量。我们最终关心的稳定对象，是相位为常数的路。这是定义 Fukaya 范畴中稳定对象时所用的特殊拉格朗日条件的一个类比。与一个稳定性条件相伴的中心荷函数，正是本文通篇一直在用的那个：$Z_\gamma = \frac{1}{\pi}\oint_\gamma\lambda$。

若有一对相交的路 $\gamma_1$、$\gamma_2$，使它们在交点附近的相位满足 $\vartheta_2<\vartheta_1<\vartheta_2+\pi$，我们就可以把交点磨光，定义一条新路 $\gamma_1\#\gamma_2$。我们预期，适当一般的 $\gamma$ 容许这样的分解
\begin{equation}
\gamma = \gamma_1 \# \cdots \# \gamma_n,
\end{equation}
其中每个 $\gamma_i$ 都是稳定的，且各组成项 $\gamma_i$ 的中心荷之相位单调递减。这类分解的存在性与唯一性，是 \cite{ks1} 处理穿墙公式的进路的一个本质前提。

获得这类分解的一种途径已在 \cite{thomas-yau} 中借助平均曲率流的一个变体描述过。本文的技术提示了另一种可能。这就是：如我们在第 6.3 节所述，% TODO 合稿改 \ref{subsec:GlobalBehaviorWKB}
对任意相位 $\vartheta$，都有 $C$ 的一个相应胞腔分解，其边界是分离性 WKB 曲线。当我们让 $\vartheta$ 顺时针转过一般值时，这些胞腔连续变化；我们同样可以连续地形变 $\gamma$，使它和分离性 WKB 曲线的关联关系保持不变。在某个临界值 $\vartheta=\vartheta_c$（此处出现一个 BPS 超多重态）处，胞腔分解的拓扑改变：其中一个胞腔塌缩。这个塌缩的胞腔可以「俘获」$\gamma$ 的一段，如图 84 所示。% TODO 图号待合稿改 \ref{fig:collapsing-cell}
这就把 $\gamma$ 自然地分裂为 $\gamma_1$ 与 $\gamma'$：$\gamma_1$ 是被俘获的那一段（它是相位为 $\vartheta_c$ 的稳定对象），$\gamma'$ 是其余部分（可能不连通）。如此继续，$\gamma'$ 将进一步分解为对象 $\gamma_i$。注意，$\gamma_i$ 自然而然地按相位有序出现，因为我们是在沿确定的方向旋转 $\vartheta$！这样我们就从几何上确定了一个不稳定对象的「Harder–Narasimhan 滤链」。

\insfig{5.6in}{collapsing-cell}{当 $\vartheta$ 变化时，WKB 叶层化所定义的某个胞腔（阴影）塌缩，俘获 $\gamma$ 在 $C$ 上投影的一段。这给出 $\gamma$ 到 $\gamma_1$ 与 $\gamma'$ 的一个自然分裂。}

如我们在第 3.3 节所评论，% TODO 合稿改 \ref{sec:iib}
人们其实预期：本文所考虑的理论可以直接由某个适当的非紧 Calabi–Yau 三维流形上的 Type IIB 弦论实现。可以设想，我们这里描述的范畴应当等同于该三维流形的 Fukaya 范畴的某个适当版本。

最后我们指出，以上所有考虑都应当自然地推广到 $K>2$ 的情形；那里不仅要考虑 BPS 弦，还要考虑 BPS 弦网。关键问题当然是确定那个结构：它推广三角剖分 $T_\WKB(\vartheta,\lambda)$、控制 WKB 近似，并编码该理论的谱生成元。自本文第一个预印本版本出现以来，我们在这个方向上已有一些进展，将出现在 \cite{gmn5-to-appear} 中。

\section{$R\to\infty$ 极限}\label{sec:R-To-Infinity}

如我们在第 2 节所回顾的，% TODO 合稿改 \ref{sec:review}
\cite{Gaiotto:2008cd} 的预期是：当 $R\to\infty$ 时应当有
\begin{equation}\label{eq:RtoInfty2}
\CX_{\gamma}^\vartheta \sim \exp\left(\pi R\zeta^{-1}Z_\gamma + i\theta_\gamma + \pi R\zeta\bar Z_\gamma\right)\left(1+\CO(e^{-\mathrm{const.}\,R})\right)
\end{equation}
只要 $\zeta$ 位于以 $e^{i\vartheta}$ 为中心的半平面 $\IH_\vartheta$ 内。我们想验证：对于我们所定义的 $\CX_\gamma^\vartheta$，\eqref{eq:RtoInfty2} 确实成立。

我们把 $\CX_\gamma^\vartheta$ 定义为 Hitchin 模空间 $\CM$ 上的函数，因此为了哪怕让方程 \eqref{eq:RtoInfty2} 有意义，我们也必须在不同 $R$ 的模空间 $\CM$ 之间作某种等同。为此我们用到这样一个事实：$\CM$ 总可以典范地等同于（与 $R$ 无关的）Higgs 丛模空间，如第 4.2.2 节所指出。% TODO 合稿改 \ref{sec:higgs}
而且，在这个模空间上还有角坐标 $\theta_\gamma$ 的一个自然定义，下面就会看到。

在让 \eqref{eq:RtoInfty2} 有意义之后，我们证明它的策略是：对相应于任意固定 $(u,\theta)$ 的 Hitchin 方程解，给出一个显式但\emph{近似}的描述。这个近似解在 $R\to\infty$ 极限下非常接近真解；事实上，通过研究线性化扰动，我们将论证它与精确解 $(A,\varphi)$ 只相差一些在 $R\to\infty$ 时指数式压低的修正。另一方面，这个近似解沿 WKB 三角剖分的边恰好是对角的，因而构成大 $R$ 下作 WKB 分析的一个非常方便的起点。这个 WKB 分析随即给出所要的 \eqref{eq:RtoInfty2}。

在本节中，我们会比本文其他部分更多地遇到 $\CM$ 的精细结构——我们不得不稍加留意诸如依赖于规范群确切选取的有限覆盖之类的问题。故事的最简单版本出现在我们把规范群取为 $PSU(2)$ 而非 $SU(2)$ 的时候，所以为简单起见，这里只限于这一情形。相关问题在附录 A 末尾讨论。% TODO 合稿改 \ref{app:MonodromyX}

\subsection{Higgs 丛的参数化} \label{sec:higgs-param}

如上所述，为了等同不同 $R$ 值下的模空间 $\CM$，方便的做法是使用明显与 $R$ 无关的 Higgs 丛图像。我们先对 Higgs 丛的规范等价类给出更显式的描述。

设有一个 Higgs 丛 $(V,\varphi,\bar\partial)$。照例写 $\Tr\varphi^2 = 2\lambda^2$，并假定 $\lambda^2$ 只含单零点。$\lambda^2$ 确定 Hitchin 纤维化的底 $\CB$ 上的一点。现在我们要参数化环面纤维。

在 $C$ 上任选一个自旋结构 $K^\half$，并把 $V$ 与 $K^{-\half}\oplus K^{\half}$ 等同。（这里自旋结构的具体选取无关紧要，因为不同选取给出同一个 $PSL(2,\IC)$ 丛。）于是至多相差规范等价，我们可以取 $\varphi\in\operatorname{End}(K^{-\half}\oplus K^{\half})\otimes K$ 为如下形式
\begin{equation} \label{eq:phi-simple}
\varphi = \begin{pmatrix} 0 & 1 \\ \lambda^2 & 0 \end{pmatrix}.
\end{equation}
这样做之后，$\bar\partial\varphi=0$ 蕴涵 $\bar\partial$ 与 $K^{-\half}\oplus K^{\half}$ 上的标准全纯结构 $\bar\partial_0$ 只相差一个与 $\varphi$ 交换的矩阵。方便的写法是
\begin{equation} \label{eq:bp-simple}
\bar\partial = \bar\partial_0 + \begin{pmatrix} 0 & \lambda^{-1} \\ \lambda & 0 \end{pmatrix} a^{0,1},
\end{equation}
其中 $a^{0,1}$ 是谱曲线 $\Sigma$ 上的 $(0,1)$-形式，且在交换两叶下为奇（从而使 $\bar\partial$ 有良好定义）。

并非所有 $a^{0,1}$ 都给出规范不等价的 Higgs 丛：我们应当把与 \eqref{eq:phi-simple} 交换的 $\operatorname{End}(K^{-\half}\oplus K^{\half})$ 的任意截面 $g_c$ 所给出的规范变换商掉。这样的 $g_c$ 可以写成
\begin{equation}
g_c = \exp\left[i\begin{pmatrix} 0 & \lambda^{-1} \\ \lambda & 0 \end{pmatrix} f\right]
\end{equation}
其中 $f$ 是 $\Sigma$ 上的复值函数，在交换两叶下为奇。$g_c$ 必须单值，$f$ 却不必：绕一个圈 $\gamma\in H_1(\Sigma;\IZ)$ 走一圈后，它可能平移
\begin{equation} \label{eq:f-multi}
f \to f + 2\pi n_\gamma
\end{equation}
其中 $n_\gamma\in\IZ$。\footnote{这里有一处微妙之处。我们实际上考虑的是结构群为 $PSL(2,\IC)$ 的 Higgs 丛，因此 $g_c$ 只需在该群中单值。于是人们也许会预期 $n_\gamma$ 应允许取半整数。但事实上，要求 $g_c$ 在 $\lambda$ 的零点附近有良好定义，会蕴涵 $n_\gamma$ 是整数。} 这样的规范变换使
\begin{equation} \label{eq:gtf}
a^{0,1} \to a^{0,1} + i\bar\partial f.
\end{equation}
利用这一自由度，我们可以安排
\begin{equation} \label{eq:a-flatness}
d a^{0,1} - \overline{d a^{0,1}} = 0,
\end{equation}
换言之，$a^{0,1}$ 其实是一个虚的\emph{平坦} 1-形式的 $(0,1)$ 部分
\begin{equation}
\tilde a = a^{0,1} - \overline{a^{0,1}}.
\end{equation}
即使在固定 \eqref{eq:a-flatness} 之后，仍剩有一些规范自由度：若 $f$ 为实，则规范变换 \eqref{eq:gtf} 保持 \eqref{eq:a-flatness}，并使
\begin{equation}
\tilde a \to \tilde a + i\,df.
\end{equation}
因此 Hitchin 纤维化的纤维由虚的平坦奇 1-形式 $\tilde a$ 模掉这一等价关系来参数化。更显式地，我们可以用
\begin{equation}
\theta_\gamma := i\oint_\gamma\tilde a
\end{equation}
给出纤维上的坐标，其中 $\gamma\in\hat\Gamma$。像 \eqref{eq:f-multi} 那样的多值规范变换把 $\theta_\gamma$ 平移为 $\theta_\gamma\to2\pi n_\gamma+\theta_\gamma$，因此 $\theta_\gamma\in\IR/2\pi\IZ$。

这样我们就得到了 Hitchin 纤维化纤维上的角坐标 $\theta_\gamma$。也可以把它们看成某个拓扑平凡复线丛上的 $\pb$-算子模掉复阿贝尔规范变换。若考虑全部 $\pb$-算子，我们得到的是 $\CA^{0,1}/\CG_c$，它正是曲线 $\Sigma$ 的雅可比簇。由于 $a^{0,1}$ 在覆叠变换下为奇，我们实际上得到的是 Prym 子簇。

\subsection{重新表述 Hitchin 方程}

如我们在第 4.2.2 节所述，% TODO 合稿改 \ref{sec:higgs}
对我们刚刚描述的每个 Higgs 丛，都有 Hitchin 方程 \cite{MR89a:32021,MR965220,MR1040197,wnh} 的一个相应解。Higgs 丛虽不依赖 $R$，其相应的解 $(A,\varphi)$ 却当然依赖，而我们关心的是研究它在 $R\to\infty$ 极限下的行为。

从 Higgs 丛 $(V,\varphi,\bar\partial)$ 到所要的 $(A,\varphi)$ 的过程如下。首先把算子 $\bar\partial$ 与 $d+A$ 的 $(0,1)$ 部分等同。其次，在 $V$ 上固定一个 Hermitian 度量，然后把这度量下 $\bar\partial$ 的伴随的负值定义为 $d+A$ 的 $(1,0)$ 部分，使得整个联络是酉的。于是可以问：这个联络连同 $\varphi$ 是否给出 Hitchin 方程的解？（它们显然满足 $\bar\partial_A\varphi=0$，所以真正的问题是 $F+R^2[\varphi,\bar\varphi]=0$ 是否成立。）对 $V$ 上度量的某种选取——称为\emph{调和度量}——情况确实如此。

指定调和度量的一种方便方式，是给出我们在第 \ref{sec:higgs-param} 节所用基与一组酉基之间的基变换矩阵 $B$。方便的做法是取一个坐标卡 $U\subset C$，局部坐标为 $z$，并取 $K^\half$ 在 $U$ 上的一个平凡化 $(dz)^\half$。然后由
\begin{equation}
\lambda^2 = p(z)(dz)^2
\end{equation}
定义 $p(z)$，并由
\begin{equation}
\eta := \left(\frac{p}{\bar p}\right)^{1/8} = \frac{p^{1/4}}{|p|^{1/4}}
\end{equation}
定义一个（多值）函数 $\eta$。

给定度量之后，酉基的选取还有一些自由度：我们用如下要求把它固定下来——在这组基中 $\varphi$ 应当纯非对角，且其右上元为实。于是最一般的基变换矩阵取形式
\begin{equation} \label{eq:M}
B = \begin{pmatrix} |p|^{\qtr} e^{h/2} & 0 \\ 0 & |p|^{-\qtr} e^{-h/2} \end{pmatrix} \exp\left[\varphi f_c/p^{1/2}\right]
\end{equation}
其中 $h$ 是 $U$ 上的实值函数。在酉基中于是得到
\begin{equation}\label{eq:RealSect}
\varphi = \begin{pmatrix} 0 & |p|^{1/2} e^h \\ \frac{p}{|p|^{1/2}} e^{-h} & 0 \end{pmatrix}
\end{equation}
以及
\begin{equation}\label{eq:SpecGaugeField}
A = a_{\bar z}\,d\bar z\begin{pmatrix} 0 & \bar\eta^2 e^h \\ \eta^2 e^{-h} & 0 \end{pmatrix} + a_z\,dz\begin{pmatrix} 0 & \bar\eta^2 e^{-h} \\ \eta^2 e^{h} & 0 \end{pmatrix} + b\,\sigma^3,
\end{equation}
其中 $a^{0,1}=a_{\bar z}d\bar z$，我们定义了 $a_z:=-(a_{\bar z})^*$，并且
\begin{equation}\label{eq:bee}
b = (\p-\pb)\log\left(|p|^{\qtr} e^{h/2}\right) = \frac{d\eta}{\eta} + \half(\p-\pb)h.
\end{equation}
代入 Hitchin 方程后便要求 $h$ 与 $f_c$ 满足
\begin{equation}\label{eq:DeformedpSG}
\p_z\pb_{\bar z}h - \left(|p|R^2 + |a_{\bar z} + \p_{\bar z}f_c|^2\right)(e^{2h}-e^{-2h}) = 0,
\end{equation}
\begin{equation}\label{eq:DefFlat}
e^{-h}\p_z\left(e^{2h}(\p_{\bar z}f_c + a_{\bar z})\right) + e^{h}\pb_{\bar z}\left(e^{-2h}(\p_z\bar f_c + (a_{\bar z})^*)\right) = 0.
\end{equation}
（注意 \eqref{eq:DefFlat} 并不等价于它的复共轭。）

与上一节一样，我们部分地固定 $a_{\bar z}$ 的选取，即要求相应的虚 1-形式 $\tilde a$ 平坦。此外，通过对 $\tilde a$ 作规范变换，我们可以取 $f_c$ 为实。\footnote{下面我们还会进一步部分地固定 $\tilde a$，即取它在转向点的某个邻域内为零。这时我们未必能同时在该邻域内取 $f_c$ 为实。不过，我们仍可继续在远离转向点的某个邻域内取它为实。}

\subsection{平凡解} \label{sec:trivial}

一个显而易见的选择是取 $h=0$。此时 \eqref{eq:DeformedpSG} 显然被满足，而 \eqref{eq:DefFlat} 变成
\begin{equation}\label{eq:flatreduced}
2\p_z\p_{\bar z}f_c + \p_z a_{\bar z} + \pb_{\bar z}a_{\bar z}^* = 0,
\end{equation}
若取 $f_c=0$，它也被满足，因为我们已经选取 $\tilde a$ 平坦。

这个解唯一的麻烦在于：它不能导致 Hitchin 方程的正则解。确实，当 $h=f_c=0$ 时，\eqref{eq:M} 给出的基变换矩阵 $B$ 在转向点（那里 $p=0$）变为奇异。要求 $B$ 正则会强加一个边界条件，而该条件将迫使 $h$ 在每个转向点处有奇性。

下一小节我们转而讨论 Hitchin 方程的这位名副其实的正则解。我们将论证：若固定一个不是转向点的点 $\CP$，则当 $R$ 足够大时，$h(\CP)$ 与 $f_c(\CP)$ 随 $R\to\infty$ 指数式快地衰减。由此可见，我们这里考虑的「平凡」解其实与精确解指数式接近。

\subsection{正则解}

为简单起见，我们先分析 $\theta_\gamma=0$ 的情形。此时可以取 $\tilde a=0$，而显然 $f_c=0$ 将满足 \eqref{eq:DefFlat}。这使 \eqref{eq:DeformedpSG} 得到重要的简化：它变成度量为 $ds^2=|p(z)||dz|^2$ 的曲面上标量场 $h$ 的 sinh-Gordon 理论的运动方程，
\begin{equation}\label{eq:SG}
\p_z\pb_{\bar z}h - R^2|p|(e^{2h}-e^{-2h}) = 0.
\end{equation}
把 \eqref{eq:SG} 乘以 $h$ 再积分即可看出：在任何紧黎曼面上都不存在 $h\neq0$ 的非奇异解。这倒也无妨，因为我们之所以想要 $h\neq0$ 的解，唯一原因就是要处理非平凡的边界条件，下面就来描述这些条件。

首先，如上所述，在 $p$ 的每个零点 $z_a$ 附近，$B$ 的正则性要求
\begin{equation}
h \sim \log|z-z_a|^{-1/2}.
\end{equation}
我们还应当讨论奇点 $\CP_i\in C$ 附近的行为。我们对 Hitchin 方程的边界条件要求 $h\to0$。于是由 \eqref{eq:bee} 可得 $A\to\frac{d\eta}{\eta}\sigma^3$。特别地，$\CP_i$ 处关于 $A$ 的边界条件只有在 $m_i^{(3)}=0$ 时才能被满足。（这正是我们所预期的：我们已经把一切 $\theta_\gamma$ 固定为 $0$，而对 $\gamma=C_i$，$m_i^{(3)}=\theta_\gamma$。）

现在我们用微扰论论证满足我们边界条件的解存在。\footnote{从物理上看它们显然存在。这些解正是 sinh-Gordon 模型在点 $z=z_a$ 处有顶点算子源时的半经典场组态。} 设 $h_{\rm ax}$ 是一个近似解。关于真解 $h=h_{\rm ax}+\delta h$ 的微扰方程为
\begin{equation}
L(h_{\rm ax})\,\delta h = E_0(h_{\rm ax}) + \CI(\delta h),
\end{equation}
其中 $L(h_{\rm ax})$ 由
\begin{equation}\label{eq:L-Operator}
L(h_{\rm ax}) := \p_z\pb_{\bar z} - 4R^2|p|\cosh(2h_{\rm ax})
\end{equation}
给出，启动微扰展开的源项为
\begin{equation}
E_0(h_{\rm ax}) = -\left(\p\pb h_{\rm ax} - 2R^2|p|\sinh(2h_{\rm ax})\right),
\end{equation}
而 $\CI$ 是相互作用项：
\begin{equation}
\CI = 2R^2|p|\left(\sinh(2h_{\rm ax})(\cosh(2\delta h)-1) + \cosh(2h_{\rm ax})(\sinh(2\delta h)-2\delta h)\right).
\end{equation}
设 $G(x,y)$ 是算子 $L(h_{\rm ax})$ 的格林函数。把它想成欧氏空间中二维有质量标量场的算子，质量平方为 $4R^2\cosh(h_{\rm ax})>4R^2$，所在曲面的度量为 $ds^2=|p(z)||dz|^2$。（在 $p$ 的零点附近，格林函数退化为质量平方为 $\mathrm{const.}\,R^2$ 的标量场的格林函数。）因此，在远离转向点处，对度 $|p(z)||dz|^2$ 下保持固定间距的 $x,y$，当 $R\to\infty$ 时有
\begin{equation}
|G(x,y)| \sim e^{-2R d(x,y)}
\end{equation}
\emph{甚至更小}（因为质量可以变大），其中 $d(x,y)$ 是从 $x$ 到 $y$ 的测地距离。另一方面，当 $d(x,y)\to0$ 时，质量对短距离行为变得无关，而 $G(x,y)\sim-\frac{2}{\pi}\log|z(x)-z(y)|$。我们有
\begin{equation}\label{eq:FirstOrderh}
\delta h(x) = \int G(x,y)E_0(h_{\rm ax}(y))\,d^2 y.
\end{equation}

现在我们来构造一个合适的 $h_{\rm ax}$。通过一个使 $dw=2R\sqrt{p(z)}dz$ 的变量替换，可以把方程化为
\begin{equation}\label{eq:2DSG}
\p_w\p_{\bar w}h = \half\sinh(2h).
\end{equation}
鉴于我们的边界条件，自然要寻找在 $w$-坐标下绕 $z_a$ 径向对称的解。于是，在某个邻域 $D_a=\{z:|z-z_a|=\rho_a\}$ 中，我们取 $w=\int_{z_a}^z 2R\sqrt{p(z)}dz$。当然，$w$ 只是局部有良好定义，但在 $z_a$ 的一个适当邻域 $D_a$ 中，它至多差一个符号没有定义，因而 $|w|$ 有良好定义。因此，在这样的邻域中寻找只依赖于 $|w|$ 的径向对称解 $h$ 是有意义的。令 $x=2|w|$，关于 $h=\tilde h(x)$ 的方程化为
\begin{equation}
\left(\frac{d^2}{dx^2} + \frac{1}{x}\frac{d}{dx}\right)\tilde h = \half\sinh(2\tilde h).
\end{equation}
这是 Painlevé III 方程，人们在把它与大质量 Ising 模型中自旋算符和无序算符的关联函数相联系方面已作过深入研究。（例如见 \cite{Fateev:1998xb}。）已知若
\begin{equation}\label{eq:conn1}
\tilde h(x) = 2\sigma\log(8/x) - \log(\gamma(\half-\sigma)) + \CO(x^{2\pm4\sigma})
\end{equation}
在 $x\to0$ 时成立，则
\begin{equation}\label{eq:conn2}
\tilde h(x) \to \frac{2\sin(\pi\sigma)}{\pi}K_0(x)
\end{equation}
在 $x\to\infty$ 时成立 \cite{McCoy:1976cd}。针对我们的边界条件，取 $\sigma=1/6$。我们把所得的、定义在 $D_a$ 中的解记作 $h_P^{(a)}$。注意，对任意固定的 $z\in D_a$，当 $R\to\infty$ 时有 $h_P^{(a)}\sim\pi^{-1}K_0(\frac{8R}{3}|z|^{3/2}(1+\cdots))$，这是指数式小的。

现在令 $C-\amalg_a D_a=\CF$ 为转向点附近这些区域之补集。我们把它称为「fatgraph 区域」（即加厚图区域），因为我们把它看作 WKB 三角剖分的一个加厚版本。取
\begin{equation}
h_{\rm ax} = \begin{cases} h_P^{(a)} & {\rm in}\quad D_a \\ 0 & {\rm in}\quad \CF \end{cases}
\end{equation}
这个选取很方便，因为 \eqref{eq:FirstOrderh} 中的积分只从 fatgraph 区域的边界 $\amalg\partial D_a$ 上获得贡献。

圆盘 $D_a$ 的边界对 $\delta h(x)$ 的贡献可以写成
\begin{equation}
-\frac{1}{4}\int_0^{2\pi}d\theta\left(\rho_a G(x,\rho_a e^{i\theta})\p_\rho h_P^{(a)} - \rho_a\left(\frac{\p}{\p\rho}G(x,\rho e^{i\theta})\right)\Big|_{\rho=\rho_a}h_P^{(a)}\right)
\end{equation}
其中我们使用坐标 $z=z_a+\rho e^{i\theta}$。一般地，若 $x$ 在度 $|p(z)dz^2|$ 下与 $\partial D_a$ 相距 $\Delta$，则差 $|\delta h-h_{\rm appxt}|$——非常粗略地说——是 $e^{-2R\Delta}e^{-\frac{8R}{3}\rho_a^{3/2}}$ 的量级，因而在区域 $D_a$ 中比 $h_P^{(a)}$ 指数式地小。

这个论证唯一失效的地方是 $x\in\partial D_a$。在短距离下可以替换 $G(x,y)\sim-\mathrm{const.}\log|x-y|$。于是人们发现，$x\in\partial D_a$ 处的修正因此是 $h_P^{(a)}$ 的量级；这相当合理，因为真解将是 $\CC^\infty$ 的，而我们的初始近似 $h_{\rm ax}$ 不连续。所以我们预期这些修正会把不连续处抹平。在其他任何地方，修正都比 $h_P^{(a)}$ 指数式地小。

于是我们断定：我们已经建立了一个良好的近似方案，并且已经证明可以自洽地取 $h=h_{\rm ax}$。

最后，让我们考虑上述讨论在 $\theta_\gamma\neq0$ 时的修正。人们可以进一步部分地固定平坦规范场 $\tilde a$ 选取中的规范自由度，即取它在转向点附近为零。然后可以验证：\eqref{eq:DefFlat} 存在一个光滑且与我们关于 $h$ 的边界条件相容的解 $f_c$。因此，我们可以继续在转向点邻域内取同样的边界条件 $h\sim-\half\log|z-z_a|$。在远离转向点处，我们把 \eqref{eq:DeformedpSG}、\eqref{eq:DefFlat} 这一对方程的解构造为关于 $h$ 与 $f_c$ 的微扰级数。若 $h=0$，我们知道 $f_c$ 非零。若 $a=0$，我们也知道可以取 $f_c=0$。若由于 $\theta_\gamma\neq0$ 而不能把 $a$ 规范掉为零，则不能令 $f_c=0$；但从微分方程可知，在远离转向点处 $f_c$ 的导数是 $h\times a$ 或更小的量级。因此我们预期，在 fatgraph 区域中当 $R\to\infty$ 时 $f_c$ 将指数式地小。注意，非零的 $\tilde a+\pb f_c-\p\bar f_c$ 只是增大了格林函数中的质量参数，从而增强了指数压低——正是这种压低使 $h$ 小、进而使 $f_c$ 小。若能做得更好、真正解出 $f_c$ 的连接问题（把 fatgraph 区域中指数式小的解与转向点附近的适当解连接起来），那将是很好的。

\subsection{$\CX_\gamma$ 的大 $R$ 极限}

上一节的结论是：沿 WKB 三角剖分的边，第 \ref{sec:trivial} 节描述的「平凡」解等于 Hitchin 方程的精确解，只差在 $R$ 下指数式压低的修正。

现在，利用 $SU(2)$ 规范变换\footnote{这个规范变换即使提升到 $\Sigma$ 上其实也是多值的；但要注意，若把它看作 $PSU(2)$ 的元素，它在 $\Sigma$ 上就是单值的。正如本节开头所指出，我们只考虑规范群 $PSU(2)$，以此回避一些微妙之处。}
\begin{equation}\label{eq:unitarygtii}
g = \frac{1}{\sqrt{2}}\begin{pmatrix} \bar\eta & -\bar\eta \\ \eta & \eta \end{pmatrix},
\end{equation}
我们可以把这个平凡解化为对角规范，其中
\begin{equation}
\varphi = p^{1/2}\sigma^3,
\end{equation}
\begin{equation}
A = \left(\tilde a + (\pb f_c - \p\bar f_c)\right)\sigma^3.
\end{equation}
然后，通过与上面第 7.4 节计算 $\CX_\gamma$ 的 $\zeta\to0$ 渐近行为时非常相似的计算，% TODO 合稿改 \ref{sec:asymptotics}
我们发现
\begin{equation}
\CX^{\vartheta}_{\gamma} = \CX_{\gamma}^{\sf}\left(1 + \CO(e^{-cR})\right),
\end{equation}
其中
\begin{equation}
\CX_{\gamma}^{\sf} = \exp\left[\pi R\zeta^{-1}Z_\gamma + i\theta_\gamma + \pi R\zeta\overline{Z_\gamma}\right].
\end{equation}
我们猜测上述常数 $c$ 的最优值由最小周期的模给出，即 $\pi\min_\gamma|Z_\gamma|$。确实，把这些修正算出来将会极其有意思！与 \cite{Gaiotto:2008cd} 的方程 (5.14) 作比较，将使人能够提取 BPS 简并度 $\Omega$。

\subsection{实截面}

在上述讨论中，所有 $\theta_\gamma=0$（从而可以取 $\tilde a=0$、$f_c=0$）的情形分析起来特别简单。这个轨迹还有另一件有趣之处。

注意，对 \eqref{eq:RealSect} 作一个把 $(\sigma^1,\sigma^2,\sigma^3)\to(\sigma^1,\sigma^3,-\sigma^2)$ 的刚性规范变换，我们可以使 $\varphi$ 成为对称矩阵，而 \eqref{eq:SpecGaugeField} 中 $a=0$ 时的 $A$ 变成实的反对称矩阵。因此，若取 $\zeta$ 为一个相位，则 $\CA$ 无迹且为实，也就是说它取值于 $sl(2,\IR)$。由 (6.10) 特别可以推出：绕每个 $\CP_i$ 的单值本征值 $\mu_i$ 都是实的。% TODO 合稿改 \eqref{eq:monodromy-formula}
于是单值矩阵是双曲的，我们可以把装饰（单值本征截面）选为实的。由此可知 Fock–Goncharov 坐标 $\CX^T_E$ 在这个轨迹上是实的。它定义了 Hitchin 纤维化的一个特殊的「实截面」。这个实截面曾在 \cite{MR89a:32021,MR1174252} 中讨论过，也是 \cite{MR2233852} 的考虑中一个非常重要的要素。

\subsection{与 Hitchin 流的关系}

作为本节的结束，我们来验证：$\zeta=0$ 处全纯辛结构中的 Hitchin 流通过 $\theta_\gamma$ 的线性平移起作用。

由方程 (1.9) 与% TODO 合稿改 \eqref{eq:symplectic}
\begin{equation}
\varpi_\zeta = -\frac{i}{2\zeta}\omega_+ + \omega_3 - \frac{i}{2}\zeta\omega_-
\end{equation}
可知，复结构 $\zeta=0$ 中的全纯辛形式为：
\begin{equation}
\omega_+ = 2iR\int\Tr\left(\delta\varphi\wedge\delta A^{0,1}\right).
\end{equation}
现在，利用 Higgs 丛的观点，考虑一个保持 $\varphi$ 不变的线性流，同时
\begin{equation}
a^{0,1} \to a^{0,1} - it\,\omega_\gamma^{0,1}
\end{equation}
这里 $\omega_\gamma$ 是 $\Sigma$ 上 Poincaré 对偶于 $\gamma$ 的实 1-形式，且在覆叠变换下为奇。这导致线性流 $\theta_{\gamma'}\to\theta_{\gamma'}+t\langle\gamma',\gamma\rangle$。把这个向量场与 $(2,0)$-形式缩并，得到
\begin{equation}
\iota\left(\frac{\p}{\p t}\right)\omega_+ = \delta\left(4R\int_C\delta\lambda\,\omega_\gamma\right) = \delta\left(2\pi R Z_\gamma\right).
\end{equation}
由此结果我们断定，辛形式可以漂亮地写成
\begin{equation}\label{eq:simplesymplectic}
\omega_+ = -2\pi R\langle dZ,d\theta\rangle + \omega_+^{\rm horizontal},
\end{equation}
其中 $\omega_+^{\rm horizontal}$ 与 Hitchin 纤维化的竖直向量缩并为零。

方程 \eqref{eq:simplesymplectic} 是对我们关于坐标 $\CX_\gamma$ 的 $R\to\infty$ 渐近行为的断言的一个很好的检验。确实，若取辛形式 $\varpi_\zeta:=\half\langle d\log\CX,d\log\CX\rangle$ 的大 $R$ 极限\footnote{本文所用的 $\varpi(\zeta)$ 归一化与 \cite{Gaiotto:2008cd} 中的选取不同。利用 \cite{Gaiotto:2008cd} 的方程 5.16，我们得到
\begin{equation}
\{\log\CX_\gamma,\log\CX_{\gamma'}\}_\zeta = 4\pi^2R\langle\gamma,\gamma'\rangle
\end{equation}
因此 $\varpi_\zeta^{\rm here} = 4\pi^2R\,\varpi_\zeta^{\rm there}$。}，便得到：
\begin{equation}\label{eq:RtoInfty3}
\varpi_\zeta \to \half\left\langle \pi R\zeta^{-1}dZ + \pi R\zeta d\bar Z + i\,d\theta,\; \pi R\zeta^{-1}dZ + \pi R\zeta d\bar Z + i\,d\theta\right\rangle
\end{equation}
我们可以从中提取 $\zeta$ 的各次幂。二重极点的系数 $\langle dZ,dZ\rangle$ 为零，因为各个 $dZ_\gamma$ 是全纯微分的周期。\footnote{注意，如果我们像 Bridgeland 稳定性条件的理论中那样，允许 $Z$ 在 $\operatorname{Hom}(\hat\Gamma,\IC)$ 中任意变化，这就\emph{不会}发生。库仑分支 $\CB$ 局部地是某个适当的 Bridgeland 稳定性条件空间的一个拉格朗日子空间。} $\varpi_\zeta$ 在 $\zeta=0$ 处单极点的留数与 \eqref{eq:simplesymplectic} 一致。

哈密顿量 $Z_\gamma$ 与标准的「Hitchin 哈密顿量」$h_a$ 密切相关。后者是 $\CB$ 上的函数，由把二次微分 $\Tr\varphi^2$ 展开为
\begin{equation}
\Tr\varphi^2 = \sum h_a\beta_a,
\end{equation}
来定义，其中 $\beta_a$ 是在 $\CP_i$ 处有二阶极点的二次微分的一组基。对 $\IC\IP^1$ 上 $\nsing$ 个正则奇点的情形，$a=1,\dots,2\nsing-3$。另一方面，我们可以展开 $\lambda=\sum_{k=1}^{3\nsing-6}Z_{\gamma_k}\alpha_k$，其中 $\alpha_k$ 是 $\Sigma$ 上对偶于 $\hat\Gamma$ 的一组基 $\gamma_k$ 的亚纯 1-形式，在交换两叶下反变，$\sigma^*(\alpha_k)=-\alpha_k$。现在，$\lambda^2=\half\Tr\varphi^2=\sum_{k,j}Z_{\gamma_k}Z_{\gamma_j}\alpha_k\alpha_j$。但 $\alpha_k\alpha_j$ 在 $\sigma$ 下不变，因而可以在基 $\beta_a$ 中展开。这样我们就能把标准哈密顿量 $h_a$ 写成周期 $Z_{\gamma_k}$ 的二次多项式。

\section{与 \texorpdfstring{\cite{Gaiotto:2008cd}}{[1]} 的比较：微分方程与黎曼--希尔伯特问题} \label{sec:comparison}

本节中，我们想比较本文所构造的达布坐标 $\CX_\gamma^\vartheta$ 的性质，与 \cite{Gaiotto:2008cd} 中确立了其存在性的坐标 $\CX^\RH_\gamma$ 的性质。

\cite{Gaiotto:2008cd} 的 $\CX_\gamma^\RH$ 有一个重要性质：（令 $\CX$ 代表任意一个 $\CX_\gamma^\RH$ 或它们的任意函数）它们满足如下形式的一组微分方程
\begin{equation}\label{eq:conn}
\p_{u^j}\CX = \left(\frac{1}{\zeta}\CA^{(-1)}_{u^j} + \CA^{(0)}_{u^j}\right)\CX,
\end{equation}
\begin{equation}
\p_{\bar u^{\bar j}}\CX = \left(\CA^{(0)}_{\bar u^{\bar j}} + \zeta\CA^{(1)}_{\bar u^{\bar j}}\right)\CX,
\end{equation}
\begin{equation}
\Lambda\p_\Lambda\CX = \left(\frac{1}{\zeta}\CA^{(-1)}_\Lambda + \CA^{(0)}_\Lambda\right)\CX,
\end{equation}
\begin{equation}
\bar\Lambda\p_{\bar\Lambda}\CX = \left(\CA^{(0)}_{\bar\Lambda} + \zeta\CA^{(1)}_{\bar\Lambda}\right)\CX,
\end{equation}
\begin{equation}
R\p_R\CX = \left(\frac{1}{\zeta}\CA^{(-1)}_R + \CA^{(0)}_R + \zeta\CA^{(1)}_R\right)\CX,
\end{equation}
\begin{equation}\label{eq:zetaeq}
\zeta\p_\zeta\CX = \left(\frac{1}{\zeta}\CA^{(-1)}_\zeta + \CA^{(0)}_\zeta + \zeta\CA^{(1)}_\zeta\right)\CX.
\end{equation}
在 \eqref{eq:conn} 的右端，我们引入了一个 $\CB\times\IC^\times\times\IR_+$ 上的联络 $\CA$。不要把它与 $C$ 上由 (1.8) 定义的 $\CA$ 混淆！% TODO 合稿改 \eqref{eq:def-CA}
联络 $\CA$ 的各个分量是 $\CM$ 的环面纤维上的复向量场，即某一组角坐标 $\theta_a$ 的基下的微分算子，在 $u,\bar u$ 固定时取值。右端我们还显式写出了联络对 $\zeta$ 的依赖。

在写这些方程时，最自然的做法是把 $\CX$ 不只看成 $\CM$ 上的函数，而看成 $\CM$ 的一个扩充版本上的函数：那里决定奇点 $\CP_i$ 处留数的参数 $m_i$、$m_i^{(3)}$ 允许变化。于是底 $\CB$ 被扩充以包含 $m_i$，而环面纤维被扩充以包含 $m_i^{(3)}$。

注意，若某一组坐标满足这些方程，则任何由与 $(R,\zeta)$ 无关的坐标变换与它相联系的其他坐标系也满足它们。特别地，若某组 $\CX^\vartheta_\gamma$ 满足它们，或者某个三角剖分对应的 $\CX^T_E$ 满足它们，那么更常规的坐标——例如单值矩阵的迹——也满足（因为这些坐标可以写成 $\CX^T_E$ 的某些有理函数——见附录 A）。% TODO 合稿改 \ref{app:MonodromyX}

从我们的渐近分析就可以直接得到这些微分方程。这里我们只讨论最重要的例子，即控制 $\zeta$ 依赖的方程 \eqref{eq:zetaeq}。\footnote{若把 $\CM$ 看成由参数 $m_i$ 与 $m_i^{(3)}$ 扩充后的 Higgs 丛模空间，则这个方程恰好表示无穷小重标 Higgs 场 $\varphi$ 的 $\IC^\times$ 作用的无穷小生成元。} 考虑一组坐标函数 $\CX_i$ 与角 $\theta_a$，并由
\begin{equation} \label{eq:def-azeta}
\CA_\zeta = \zeta\frac{\p\CX_i}{\p\zeta}\left[\frac{\p\CX_i}{\p\theta_a}\right]^{-1}\frac{\p}{\p\theta_a}
\end{equation}
在环面纤维上定义一个向量场。由链式法则，这个定义显然既不依赖于环面纤维的具体参数化，也不依赖于坐标系 $\CX_i$ 的选取。对良好的坐标系 $\CX$，雅可比矩阵 $\frac{\p\CX}{\p\theta}$ 确实可逆，因为在 $\zeta=0,\infty$ 之外，$\CM$ 在任意复结构 $J^{(\zeta)}$ 下局部地就是环面纤维的复化。于是我们差不多是同义反复地有
\begin{equation}
\zeta\p_\zeta\CX = \CA_\zeta\CX. \label{eq:zetaeqtaut}
\end{equation}
现在，视需要确定 $\CA_\zeta$ 的哪些性质而定，合适的坐标系 $\CX$ 各有不同。为证明 $\CA_\zeta$ 在 $\zeta=0,\infty$ 之外全纯，方便的做法是：在 $\zeta$ 的每个取值处、在 $\CM$ 的每一点附近，使用一个在该点附近良好的坐标系。只要装饰给出的截面 $s_i$、$s_j$ 不在 $T$ 的某条边 $E_{ij}$ 上重合，Fock–Goncharov 坐标系 $\CX^T$ 就是适用的；此外，绕 $C$ 上各个圈的单值矩阵之迹提供了一个完全合理的整体坐标选取。为证明 $\CA_\zeta$ 在 $\zeta=0,\infty$ 处至多有一阶极点，我们考虑坐标系 $\CX^\vartheta$，同时让 $\zeta$ 在 $\IH_\vartheta$ 内趋于 $0$ 或 $\infty$。把 $\CX^\vartheta$ 的已知渐近式代入 \eqref{eq:def-azeta}，便得到关于 $\CA_\zeta$ 的所要信息。于是我们得到所要的 \eqref{eq:zetaeq}。

这组相容微分方程 \eqref{eq:conn}–\eqref{eq:zetaeq} 相当有力。在有限维的情形中，若算子 $\CA$ 是矩阵而非微分算子，这样的方程会直接等价于一个黎曼–希尔伯特问题。不过，在我们的无穷维情形中，微分方程 \eqref{eq:zetaeq} 的强度不足以保证解在 $\zeta=0,\infty$ 之外关于 $\zeta$ 全纯。（这一点从如下事实即可看出：某些 $\CX$ 的任何有理函数，只要满足 \eqref{eq:conn}–\eqref{eq:zetaeq}，也就满足 \eqref{eq:conn}–\eqref{eq:zetaeq}。）

当 $R$ 足够大时，稍加努力就可以为 WKB 分析的误差给出界，因此应当能够保证：WKB 曲线一端的那个小平坦截面，绝不会与另一端的那个小平坦截面重合；这意味着 $\CX^\vartheta_\gamma$ 在射线 $\zeta\in e^{i\vartheta}\IR_+$ 的邻域内没有极点。不过，当 $R$ 小时，我们看不出任何明确的办法来排除这种可能性。确实，WKB 三角剖分只携带关于 Higgs 场 $\varphi$ 的信息：关于规范联络 $A$ 的任何约束都只能非常间接地、通过 Hitchin 方程的解而得到。可以设想：尽管如此，也许从来不存在一个在 $T_\WKB$ 某条边的两端都很小的平坦截面 $s$；但我们既没有找到支持这一猜想的迹象，也没有找到反对它的迹象。

结果是，我们只能有把握地断言：当 $R$ 足够大时，坐标 $\CX_\gamma$——从而 Hitchin 系统模空间上的度量——可以由 \cite{Gaiotto:2008cd} 中表述的一般黎曼–希尔伯特问题，结合本文算出的谱生成元（斯托克斯矩阵）来确定。弄清楚究竟是黎曼–希尔伯特问题在小 $R$ 时干脆无解，还是解虽存在却给出错误的度量，将会很有意思。我们希望在未来一篇论文中给出一些数值检验。
